\inicializarSubseccion{Estructura b\'asica del c\'odigo}

\tab El ciclo for se usa cuando se conoce cuantas veces se repetirá un proceso, en Matlab tiene la estructura:
\begin{lstlisting}[style=Matlab-editor]
for variable = inicio:incremento:fin
    instrucciones
end
\end{lstlisting}

\tab Las partes de este son la variable, la cual lleva el conteo de iteraciones (usualmente $i$, $j$), el inicio tiene el valor inicial de la variable, el incremento indica cuanto aumento hay por iteración, normalmente de uno, el fin es el valor límite y el bloque de instrucciones es lo que se ejecutará en cada repetición.

Un ejemplo del ciclo for sería:
\begin{lstlisting}[style=Matlab-editor]
for i = 1:10
    disp(i)
end
\end{lstlisting}

\tab Por el otro lado, el ciclio while se usa cuando no se sabe con exactitud cuantas veces se va a repetir el proceso, por lo cual depende de una condición

Este ciclo tiene la estructura general:

\begin{lstlisting}[style=Matlab-editor]
while condición
    instrucciones
end
\end{lstlisting}

\tab Las partes de este son la condición que se evaluará antes de cada iteración y el bloque de instrucciones que se ejecuta mientras esta condición sea verdadera

Un ejemplo de este ciclo sería:

\begin{lstlisting}[style=Matlab-editor]
i = 1;
    while i <= 10
    disp(i)
    i = i + 1;
end
\end{lstlisting}

